\section{Kết luận}

Báo cáo được tổ chức quanh một bài toán chung về phát triển bền vững giai đoạn 2000--2023. Bốn phần kinh tế, giáo dục, sức khỏe--dân số và môi trường không phải là bốn bài toán độc lập, mà là bốn nhóm bằng chứng để đánh giá mức độ phát triển của các quốc gia và khu vực từ nhiều khía cạnh. Kết quả kinh tế hiện cho thấy sự phân hóa rõ rệt về GDP bình quân đầu người, khác biệt tăng trưởng giữa các khu vực, và xu hướng nghịch giữa GDP bình quân đầu người với tỷ lệ nghèo trong nhóm quốc gia có đủ dữ liệu. Các phần còn lại sẽ được hoàn thiện bằng ảnh dashboard, nhận xét định lượng và thảo luận giới hạn dữ liệu theo cùng một chuẩn trình bày.

Các kết luận trong báo cáo cần được hiểu như kết quả mô tả dựa trên dữ liệu WDI và dashboard Tableau, không phải bằng chứng nhân quả. Một số chỉ số, đặc biệt là tỷ lệ nghèo, có mức độ thiếu dữ liệu đáng kể theo năm và quốc gia. Vì vậy, phần thảo luận cuối cùng của nhóm cần nêu rõ phạm vi dữ liệu được sử dụng, tránh khái quát vượt quá các quan sát có sẵn, và nhấn mạnh rằng phát triển bền vững chỉ có thể được đánh giá đầy đủ khi đặt đồng thời các kết quả kinh tế, xã hội và môi trường cạnh nhau.

\section*{Tài liệu tham khảo}

\begin{thebibliography}{9}

\bibitem{worldbankwdi}
World Bank. \textit{World Development Indicators}. Washington, DC: World Bank.

\bibitem{course}
Khoa Công nghệ Thông tin, Trường Đại học Khoa học Tự nhiên, ĐHQG-HCM. \textit{Tài liệu hướng dẫn môn Trực quan hóa dữ liệu}.

\bibitem{tableau}
Tableau Software. \textit{Tableau Documentation}.

\end{thebibliography}
